\documentclass[10pt,a4paper]{amsart}
\usepackage[margin=19mm]{geometry}
\usepackage{amsmath,amssymb,mathtools}
\usepackage[scr=rsfs,cal=euler]{mathalfa}
\usepackage{xfrac}
\usepackage[unicode,hidelinks,bookmarks=false]{hyperref}
\newcommand{\ie}{i.e.\ }
\newcommand{\cf}{cf.\ }
\newcommand{\resp}{respectively\ }
\newcommand{\Subsection}{\subsection}
\providecommand{\nobreakdash}{\nobreak\hbox{-}}
\setlength{\emergencystretch}{2em}
\multlinegap0pt
\usepackage{amssymb}
\usepackage{enumerate}
\usepackage{xcolor}
\newtheorem{proposition}{Proposition}
\newcommand{\Q}{{\mathbb Q}}
\newcommand{\T}{{\mathbb T}}
\newcommand{\Z}{{\mathbb Z}}
\newcommand{\brac}[1]{\mathopen{}\left( #1 \mathclose{}\right)}
\title{Limiting behavior of 2-step nilpotent ergodic averages with polynomial iterates}
\author{Andreas Koutsogiannis, Borys Kuca, Wenbo Sun}
\date{Source: arXiv:2607.29368v2 (CC BY 4.0)}
\begin{document}
\maketitle

\noindent\emph{Source and licence.} Limiting behavior of 2-step nilpotent ergodic averages with polynomial iterates. Version arXiv:2607.29368v2 (CC BY 4.0), distributed by arXiv under the Creative Commons Attribution 4.0 International licence, http://creativecommons.org/licenses/by/4.0/, as recorded in the arXiv licence metadata for this exact version and shown on the abstract page https://arxiv.org/abs/2607.29368v2. Full licence notice: \texttt{licenses/arxiv-2607.29368-CC-BY-4.0.txt}.\par\smallskip

\noindent\textit{Editorial scope.} Counted unit: the article's proposition affine (source0.tex lines 2749--2772). The article's complete original proof of it is delivered with it (source0.tex lines 2773--2832, one \texttt{proof} environment of 60 lines). No mathematics is written, repaired or re-derived: the statement and the proof are the article's own lines, in the article's own notation, and no retained line is substituted or re-flowed.  No further source-local context is needed: every cross-reference of every retained line resolves inside the two delivered spans.  Nothing was omitted from any retained span, so the package carries no omission record.  No retained line contains a citation, so the wrapper carries no bibliography and no bibliography entry.  No retained line contains a figure environment, an image-inclusion command or drawing code, so no image is delivered and the package declares no asset dependency.  Editorial formatting only, in addition: an AMS \texttt{amsart} preamble that restates the article's own declarations and macros used by the retained lines, namely the environments \texttt{proposition} on separate counters, together with the standard packages those lines need.  The retained environments print in this document as Proposition 1 in order of appearance, whereas the article prints the same items with the numbering of its own full document.  The complete native source of the article as distributed in the arXiv e-print for this exact version is bundled unchanged as \texttt{sources/206478/source0.tex}.
\begin{proposition}\label{P: equidistribution for unipotent affine}
    For $1\leq i\leq s$, let $m_i\in\T$ and $Y_i := \{0\}^{m_{i-1}}\times \T^{m_{i}+ \cdots + m_s}$ (with $m_0:=0$). Set also $m:=m_1 + \cdots + m_s$.
    Let $p_1, \ldots, p_\ell\in\Z[n]$ be independent polynomials, and let $S_1, \ldots, S_\ell:\T^m\to\T^m$ be ergodic unipotent affine transformations of the form
    \begin{align*}
        S_jy = \alpha_j + (I+A_j + \cdots + A_j^{s-1})y
    \end{align*}
    for   $\alpha_{j}\in\T^{m}$ and some continuous group homomorphism $A_j:\T^m\to\T^m$ satisfying
    \begin{align}\label{E: shifting property of A}
     A_j(Y_i)\subseteq Y_{i+1}   
    \end{align}
    %$A_j(Y_i)\subseteq Y_{i+1}$ 
    for every $1\leq i\leq s$.    
    % \begin{align*}
    %     %S_j(x_1, \ldots, x_s) 
    %     S_j(x) := (x_1 + \alpha_{j,1}, x_2 + A_{j,2,1}x_1 + \alpha_{j,2}, \ldots, x_s + A_{j,s,1}x_1 + \cdots + A_{j,s,s-1}x_{s-1} + \alpha_{j,s})
    % \end{align*}
    % (where $x_i\in\T^{m_i}$ for each $1\leq i\leq s$)
    % for some integer $m_j\times m_i$-matrices $A_{i,j}$ and elements $\alpha_i\in \T^{m_i}$.
    Then for almost every $y\in\T^m$, the orbit
    \begin{align}%\label{E: orbit}
        (S_1^{p_1(n)}y, \ldots, S_\ell^{p_\ell(n)}y)_{n\in\Z}
    \end{align}
    is dense in $\T^{m\ell}$.
\end{proposition}

\begin{proof}
    Let $R_j:=I+ A_j + \cdots + A_j^{s-1}$. Then $R_j^n=\sum_{i=0}^{s-1}\binom{n+i-1}{i}A_j^i$ (for $n\geq 1)$ and 
    \begin{align*}
        S_j^n y &= R_j^n y + \sum_{i=0}^{n-1}R^i\alpha_j %\sum_{i=0}^n \brac{\binom{n}{i}R_j^i y + \binom{n}{i+1}R_j^i\alpha_j}\\
        = \sum_{i=0}^{s-1}\binom{n+i-1}{i}A_j^iy + n\alpha_j + \sum_{i=1}^{s-1}\binom{n+i-1}{i+1}A_j^i\alpha_j.
    \end{align*}
    %noting that $A_j^s = 0$, so all the terms past that in fact vanish.
    

    Choose $y=(y_1, \ldots, y_s)$ (with $y_{j}\in\T^{m_{j}}$ for $1\leq j\leq s$) so that the coordinates of $y$ are all $\Q$-independent from $\alpha_1, \ldots, \alpha_s, 1$, in that
    \begin{align*}
        k\cdot y + l_1 \alpha_1 + \cdots + l_s \alpha_s + l_{s+1} = 0
    \end{align*}
    for some $k\in\Z^m$ and $l_1, \ldots, l_{s+1}\in\Z$ only if $k = 0$.
    Such points form a full measure set. Suppose that
    \begin{align*}%\label{E: linear combo}
        k_1\cdot S_1^{p_1(n)}y + \cdots + k_\ell \cdot S_\ell^{p_\ell(n)}y = 0
    \end{align*}
    for some $k_1, \ldots, k_\ell\in\Z^m$. Given the formula for $S_j^n y$'s above, this becomes
    \begin{multline}\label{E: linear combo 2}
        \sum_{j=1}^\ell \left(\sum_{i=0}^{s-1}\binom{p_j(n)+i-1}{i}(k_j\cdot A_j^i)y\right.\\
        \left.+ p_j(n)(k_j\cdot\alpha_j) + \sum_{i=1}^{s-1}\binom{p_j(n)+i-1}{i+1}(k_j\cdot A_j^i)\alpha_j\right) = 0.
    \end{multline}
    Since all the coordinates are polynomial in $n$, by Weyl's equidistribution theorem it suffices to show that $k_1 = \cdots = k_\ell = 0$.

    Because of the $\Q$-independence assumption on the coordinates, the contribution of each coordinate of $y$ to \eqref{E: linear combo 2} is 0. We analyze one-by-one the contribution of coordinates $y_{s-1}, \ldots, y_1$. Since $A_j^2(Y_{s-1})\subseteq A_j(Y_s) = \{0\}$, the coordinate $y_{s-1}$ does not see the contribution of the $A_j^2, A_j^3,\ldots$
    Therefore, this coordinate contributes exactly
    \begin{align}\label{E: contribution of s-1}
        %(p_1(n) k_{1}\cdot A_{1} + \cdots + p_\ell(n) k_{\ell}\cdot A_{\ell})y_{s-1} {
        (p_1(n) k_{1}\cdot A_{1} + \cdots + p_\ell(n) k_{\ell}\cdot A_{\ell})(0,\ldots, 0, y_{s-1}, 0)
    \end{align}
    plus terms independent of $n$ that are of no interest to us. Because of the assumption on the rational independence of coordinates, \eqref{E: linear combo 2} and \eqref{E: contribution of s-1}, and the fact that $A_{j}(Y_{s})=\{0\}$, we have
    %  \begin{align} 
    %     %(p_1(n) k_{1}\cdot A_{1} + \cdots + p_\ell(n) k_{\ell}\cdot A_{\ell})y_{s-1} {
    %     (p_1(n) k_{1}\cdot A_{1} + \cdots + p_\ell(n) k_{\ell}\cdot A_{\ell})(0,\ldots, 0, y_{s-1}, y_{s})=0
    % \end{align}
    \begin{align*}
        (p_1(n) k_{1}\cdot A_{1} + \cdots + p_\ell(n) k_{\ell}\cdot A_{\ell})({Y_{s-1}}) = \{0\}
    \end{align*}
    %\eqref{E: contribution of s-1} 
     for all $n$. By the independence of polynomials, it follows that $(k_j\cdot A_j)({Y_{s-1}}) = 0$ for every $1\leq j\leq \ell$. Hence also
    \begin{align*}
        (k_j\cdot A_j^{s-1})(Y) \subseteq (k_j\cdot A_j)(Y_{s-1}) = \{0\}
    \end{align*}
    due to the property \eqref{E: shifting property of A} of the maps $A_j$.
    The consequence of this is that we can lower the upper indices in both sums in \eqref{E: linear combo 2} from $s-1$ to $s-2$. 

    We now repeat the same procedure, this time analyzing the contribution of $y_{s-2}$ to \eqref{E: linear combo 2}. Because of the previous case, the contribution of $y_{s-2}$ is the same as \eqref{E: contribution of s-1}, with $(0, \ldots, 0, y_{s-1},0)$ replaced by $(0,\ldots, 0, y_{s-2},0,0)$, plus terms constant in $n$. Arguing exactly as in the previous case, we conclude that
        \begin{align*}
        (k_j\cdot A_j^{s-2})(Y) \subseteq (k_j\cdot A_j)(Y_{s-2}) = \{0\}.
    \end{align*}
    
    Iterating this procedure, we eventually deduce that $k_j\cdot A_j = 0$ for every $1\leq j\leq \ell$. With that, \eqref{E: contribution of s-1} reduces to
    \begin{align*}
        \sum_{j=1}^\ell p_j(n)(k_j\cdot\alpha_j) = 0.
    \end{align*}
    The linear independence of $p_j$'s implies that $k_j\cdot \alpha_j = 0$ for every $1\leq j\leq \ell$.

    We have thus established that $k_j \cdot A_j = k_j \cdot \alpha_j = 0$ for every $1\leq j\leq \ell$. In particular, $k_j\cdot (S_j^n y) = k_j\cdot y$ for every $n\in\Z$. The unique ergodicity of $S_j$ implies that $k_j = 0$, and so the equidistribution of the orbit follows from Weyl's equidistribution theorem. 
 \end{proof}

\end{document}
